\section*{Bản tin tiếng Việt}
\textbf{Tên đề tài:} \ProjectTitle\\
\textbf{Mã số:} \ProjectCode\\
\textbf{Chủ nhiệm:} \PrincipalName\\
\textbf{Đơn vị:} \PrincipalUnit\\
\textbf{Người hướng dẫn:} \SupervisorName\ (\SupervisorEmail)\\
\textbf{Thời gian đăng ký:} \ProjectPeriod

\subsection*{Mục tiêu}
Xây dựng hệ thống hỗ trợ tiếp nhận, phân tích và gom nhóm báo cáo cứu hộ bão lũ trong điều kiện kết nối không ổn định, kết hợp AI tại thiết bị, truyền thông tin gọn và dashboard điều phối.
\subsection*{Phương pháp}
Huấn luyện và đánh giá phân loại ảnh trên split cố định; đối chiếu các định dạng, lượng tử hóa và cắt tỉa; xây dựng pipeline đồ thị/điểm ưu tiên từ dữ liệu quan sát; đánh giá phân cụm, xếp hạng và lịch điều phối bằng các bộ dữ liệu mô phỏng riêng; phát triển ứng dụng Flutter và backend FastAPI với lưu cục bộ, retry và hợp đồng đồng bộ.
\subsection*{Kết quả và sản phẩm}
Mô hình ảnh FP32 đạt Accuracy \ImageAccuracy\% và macro-F1 \ImageMacroFOne\% trên \ImageTest{} ảnh test; ONNX/ExecuTorch giữ cùng nhãn top-1 trên tập này. PTQ giảm dung lượng nhưng giảm chất lượng; QAT hiện có kết quả không thuận lợi. Gói LR khẩn cấp dùng 32 trạng thái tổng hợp có cấu trúc, không thay mô hình hiểu tin nhắn tiếng Việt.

Repo có ứng dụng, backend, dashboard bản đồ, mô hình/asset, dữ liệu và bản thảo nghiên cứu. Pipeline có chặn cạnh và bất biến bản sao có điều kiện; RQ1--RQ3 chưa xác lập ưu thế chung hay lợi ích cứu hộ thực tế. Kết quả mô phỏng và latency CPU được ghi riêng với tình trạng kiểm chứng thiết bị.
Summary PTQ/QAT có chênh lệch với CSV dự đoán; cần rà soát trước khi chốt metric nén.
\subsection*{Khả năng ứng dụng và nội dung tiếp tục}
Nguyên mẫu phục vụ nghiên cứu và trình diễn tiếp nhận--điều phối. Cần bổ sung dữ liệu thực địa, chuyên gia, mô hình văn bản, thống nhất trường khẩn cấp app/backend, phép đo Android, mạng thật và SMS gateway.

\DraftNote{Bản tin đã có nội dung sơ bộ; cần chuyển sang mẫu chính thức và xác nhận minh chứng trước khi nộp.}
